\section{System Design}
\label{sec:design}

\subsection{Design Goals}

Our design addresses five goals derived from the problems identified in Section 1.2:

\begin{enumerate}
    \item \textbf{Strong Isolation}: Each agent runtime must be isolated from others; a failure in one must not affect others.
    \item \textbf{Independent Scalability}: Agent runtimes should be independently scalable based on demand.
    \item \textbf{Standardized Orchestration}: An Orchestrator should be able to coordinate heterogeneous agent runtimes without knowing their internal implementation.
    \item \textbf{Security by Design}: The architecture should mitigate the security risks identified by CNCERT/CERT.
    \item \textbf{Agent-Native Abstractions}: The platform should provide abstractions (Skills, lifecycle management) tailored for agents.
\end{enumerate}

\subsection{Architecture Overview}

We design a three-layer hierarchical architecture, as shown in Figure~\ref{fig:architecture}.

\begin{figure}[htbp]
    \centering
    \includegraphics[width=0.9\textwidth]{figures/architecture.pdf}
    \caption{Three-layer hierarchical architecture: Orchestrator, Agent Runtime, and Container Platform.}
    \label{fig:architecture}
\end{figure}

\textbf{Orchestrator Layer}: The Orchestrator is itself an agent responsible for task decomposition, planning, and coordinating specialized agent runtimes through a standardized Skill interface.

\textbf{Agent Runtime Layer}: Each agent runtime (OpenClaw, Claude Code, etc.) runs in its own container. Agents do not directly invoke each other; instead, they receive instructions from the Orchestrator and execute within their containerized environment.

\textbf{Container Platform Layer}: Kubernetes manages the lifecycle, scheduling, networking, and security of each agent container.

\subsection{Layer Responsibilities}

\subsubsection{Orchestrator Layer}
The Orchestrator performs high-level task coordination:
\begin{itemize}
    \item \textbf{Task Decomposition}: Breaks complex user requests into subtasks
    \item \textbf{Skill Discovery}: Identifies which Skills are needed for each subtask
    \item \textbf{Execution Scheduling}: Invokes Skills in the appropriate agent containers
    \item \textbf{Result Aggregation}: Collects outputs from multiple agents
\end{itemize}

Critically, the Orchestrator does \textbf{not} execute tasks directly. Instead, it \textbf{dispatches} work to agent containers through Skills.

\subsubsection{Agent Runtime Layer}
Each agent container provides a complete agent runtime environment:
\begin{itemize}
    \item \textbf{OpenClaw Container}: OpenClaw runtime with memory, tools, and Skills
    \item \textbf{Claude Code Container}: Claude Code runtime in a separate container
    \item \textbf{Custom Agent Container}: Supports arbitrary agent runtimes
\end{itemize}

Agents are \textbf{specialized} and \textbf{stateless} regarding orchestration.

\subsubsection{Container Platform Layer}
Kubernetes provides:
\begin{itemize}
    \item \textbf{Lifecycle Management}: Create, start, stop, restart, and delete containers
    \item \textbf{Resource Isolation}: CPU, memory, GPU limits per container
    \item \textbf{Network Isolation}: Each container has its own network namespace
    \item \textbf{Namespace Isolation}: Each user gets a dedicated K8s namespace
    \item \textbf{Persistent Storage}: Volume mounts for agent memory and configuration
\end{itemize}

\subsection{Agent Lifecycle}

Each agent instance follows a defined lifecycle:

\begin{figure}[htbp]
    \centering
    \begin{tikzpicture}[node distance=1.5cm, auto]
        \node[draw, rounded rectangle] (create) {CREATE};
        \node[draw, rounded rectangle, right=of create] (init) {INITIALIZE};
        \node[draw, rounded rectangle, right=of init] (running) {RUNNING};
        \node[draw, rounded rectangle, right=of running] (stopped) {STOPPED};
        \node[draw, rounded rectangle, right=of stopped] (delete) {DELETE};
        
        \draw[->] (create) -- (init);
        \draw[->] (init) -- (running);
        \draw[->] (running) -- (stopped);
        \draw[->] (stopped) -- (delete);
        
        \draw[->, dashed] (running.north) .. controls +(0,0.5) and +(0,0.5) .. (create.north);
    \end{tikzpicture}
    \caption{Agent lifecycle state machine.}
    \label{fig:lifecycle}
\end{figure}

\begin{enumerate}
    \item \textbf{CREATE}: Kubernetes creates the Pod and container, allocates resources.
    \item \textbf{INITIALIZE}: Agent-specific initialization runs (via Skill).
    \item \textbf{RUNNING}: The agent is ready to receive and execute tasks.
    \item \textbf{STOPPED}: The container is stopped but resources are retained.
    \item \textbf{DELETE}: Resources are released.
\end{enumerate}

\subsection{Skill Framework}

\subsubsection{What is a Skill?}
A \textbf{Skill} is a standardized capability unit that exposes an agent's functionality to the Orchestrator:

\begin{itemize}
    \item \textbf{Name}: A unique identifier (e.g., \texttt{openclaw-manager}, \texttt{hai-k8s-container})
    \item \textbf{Description}: Human- and machine-readable description of capabilities
    \item \textbf{Interface}: Input parameters and output format
    \item \textbf{Execution Endpoint}: Where and how the Skill is invoked
\end{itemize}

Unlike traditional tools (which are local extensions), Skills are \textbf{transportable} and \textbf{orchestrator-callable} across containers.

\subsubsection{Skill Interface}
Skills are invoked through a standardized API:

\begin{verbatim}
POST /api/skills/containers/{container_id}/exec
Headers:
  X-Admin-API-Key: <admin_key>
  Authorization: Bearer <user_jwt>
Body:
  {"command": "<skill-specific-command>", "timeout": <seconds>}
\end{verbatim}

Two-factor authentication ensures that Skills can only be invoked by authorized callers on containers owned by the authenticated user.

\subsubsection{hai-k8s-container Skill}
Provides low-level container management:
\begin{itemize}
    \item \texttt{list\_containers()}: List user's containers
    \item \texttt{get\_container()}: Get container details
    \item \texttt{exec\_in\_container()}: Execute command inside container
    \item \texttt{delete\_container()}: Stop and delete container
\end{itemize}

\subsubsection{openclaw-manager Skill}
Provides OpenClaw-specific operations:
\begin{itemize}
    \item \texttt{init\_openclaw()}: Full 4-step initialization
    \item \texttt{config\_models()}: Configure model providers
    \item \texttt{get\_status()}: Check OpenClaw status
\end{itemize}

\subsection{Security Model}

Our security architecture addresses the threats identified by CNCERT/CERT:

\begin{itemize}
    \item \textbf{Container Isolation}: Each agent runs in an isolated K8s Pod; compromise of one container does not grant access to others.
    \item \textbf{Two-Layer Authentication}: Admin API Key + User JWT prevents unauthorized Skill invocation.
    \item \textbf{User Namespace Isolation}: Each user operates within a dedicated K8s namespace.
    \item \textbf{Capability Restrictions}: Non-root user, restricted Linux capabilities, read-only root filesystem.
    \item \textbf{Audit Logging}: All Skill invocations and exec operations are logged.
\end{itemize}

\subsection{Comparison with Single-Runtime Architecture}

Table~\ref{tab:comparison} summarizes the key differences.

\begin{table}[htbp]
    \centering
    \caption{Architectural comparison: Single-Runtime vs Multi-Runtime.}
    \label{tab:comparison}
    \begin{tabular}{@{}lcc@{}}
        \toprule
        \textbf{Dimension} & \textbf{Single-Runtime} & \textbf{Multi-Runtime (Ours)} \\
        \midrule
        Agent Isolation & Shared memory & Complete isolation \\
        Concurrency & Pseudo-parallel & True parallel \\
        Scalability & Modify process code & Register new container \\
        Resource Control & Process limits & Kubernetes quotas \\
        Security Boundary & Process & Container/Namespace \\
        Agent Heterogeneity & Limited & Unlimited \\
        Cold Start & Process fork & Pod creation \\
        Orchestration & Function calls & Standardized Skill API \\
        \bottomrule
    \end{tabular}
\end{table}
